% !TEX program = xelatex
% Báo cáo tiến độ — tái hiện bài báo ILCHO
% Biên dịch bằng XeLaTeX (cần font hỗ trợ tiếng Việt, vd. Times New Roman có sẵn trên Windows).
\documentclass[12pt,a4paper]{article}

\usepackage{fontspec}
\setmainfont{Times New Roman}

\usepackage[a4paper,margin=2.5cm]{geometry}
\usepackage{booktabs}
\usepackage{array}
\usepackage{enumitem}
\usepackage{caption}
\renewcommand{\tablename}{Bảng}
\usepackage{titlesec}
\usepackage{setspace}
\onehalfspacing

\titleformat{\section}{\large\bfseries}{\thesection.}{0.5em}{}
\titleformat{\subsection}{\normalsize\bfseries}{\thesubsection.}{0.5em}{}

\setlength{\parindent}{0pt}
\setlength{\parskip}{6pt}

\title{\textbf{BÁO CÁO THỬ NGHIỆM}\\[4pt]
\large Tái hiện bài báo ``Intelligent Handover Scheme for Improved 6G NTN\\
LEO Satellite Network Performance'' (ILCHO)}
\author{Sinh viên: ........................................\\
Giảng viên hướng dẫn: Cô ........................................}
\date{}

\begin{document}
\maketitle
\vspace{-2em}

\section{Bài toán}

Bài báo gốc: M. Choi, M. Park, J. Kim, J.-M. Chung, \emph{``Intelligent Handover Scheme for
Improved 6G NTN LEO Satellite Network Performance,''} IEEE Transactions on Mobile Computing,
vol. 25, no. 5, 05/2026, tr. 6863--6880.

Bài báo đề xuất \textbf{ILCHO}: kết hợp \textbf{Conditional Handover (CHO)} của 3GPP với
\textbf{học tăng cường đa tác tử (QMIX)} để chọn vệ tinh đích cho từng thiết bị người dùng
(UE) trong một chòm vệ tinh LEO mật độ cao (kiểu Starlink). Bài toán đặt ra: khi số vệ tinh
và số UE lớn, việc chuyển giao (handover) liên tục là bắt buộc (vài chục giây một lần), các
phương pháp truyền thống (chọn theo khoảng cách gần nhất, theo thời gian nhìn thấy lâu nhất)
gây ra tỉ lệ chuyển giao thất bại (HOF) cao khi tải hệ thống tăng. Mục tiêu của ILCHO là dùng
QMIX để chọn vệ tinh đích thông minh hơn, giữ HOF thấp và thông lượng (throughput) ổn định
ngay cả khi số UE tăng cao.

Bài báo \textbf{không công khai mã nguồn}, và \textbf{không công bố một số tham số quan
trọng} (tần số sóng mang chính xác, ngưỡng thực thi chuyển giao, các hằng số trong hàm
thưởng, hệ số pha của chòm vệ tinh Walker-Delta). Nhiệm vụ đặt ra là \textbf{tái hiện lại
toàn bộ phương pháp từ đầu} chỉ dựa trên nội dung bài báo, đối chiếu số liệu tái hiện được
với số liệu bài báo công bố, và báo cáo trung thực phần nào khớp / không khớp.

\section{Mục tiêu tái hiện}

\begin{itemize}[leftmargin=1.2em]
  \item Dựng lại toàn bộ mô hình: hình học quỹ đạo vệ tinh, mô hình kênh truyền (3GPP TR
        38.811/38.821), môi trường mô phỏng CHO đa tác tử, bộ điều khiển QMIX, và 4 phương
        pháp so sánh (baseline) mà bài báo dùng.
  \item Huấn luyện và đánh giá ở quy mô càng gần bài báo càng tốt (số UE, số vệ tinh, số
        episode huấn luyện).
  \item Đối chiếu các tuyên bố chính của bài báo (số chuyển giao, số chuyển giao thất bại,
        thông lượng, độ công bằng) với kết quả tái hiện được.
  \item Báo cáo trung thực: phần nào tái hiện được, phần nào không, và vì sao (do quy mô tính
        toán, hay do khác biệt thật về phương pháp/tham số ẩn).
\end{itemize}

\section{Công việc đã thực hiện}

\subsection{Dựng phương pháp từ bài báo}

Đã dựng toàn bộ mã nguồn (không có mã nguồn gốc để tham khảo) gồm: hình học chòm vệ tinh
Walker-Delta, mô hình liên kết vô tuyến theo chuẩn 3GPP, môi trường mô phỏng CHO đa tác tử,
bộ học QMIX (mạng nơ-ron hồi quy + mạng trộn giá trị đơn điệu), và 4 phương pháp so sánh:
MD-CHO (chọn theo khoảng cách gần nhất), MVT-CHO (chọn theo thời gian nhìn thấy lâu nhất),
HSNF (gán tham lam tối đa hoá thông lượng), LBSH (học tăng cường đơn tác tử + cân bằng tải).

\subsection{Ba lượt tái hiện, tăng dần quy mô}

Do giới hạn phần cứng ban đầu (chỉ có CPU), việc tái hiện được chia làm nhiều lượt, mỗi lượt
tiến gần hơn tới đúng quy mô bài báo công bố (Bảng~\ref{tab:rounds}). Lượt 3 thực hiện sau khi
đổi sang máy có GPU (RTX 3060 Ti), cho phép huấn luyện đúng số tác nhân (agent) tối đa và đúng
độ dài một lượt mô phỏng (600 giây) mà bài báo dùng.

\begin{table}[h]
\centering
\caption{Ba lượt tái hiện, tiến dần tới đúng quy mô bài báo}
\label{tab:rounds}
\begin{tabular}{lccc}
\toprule
 & \textbf{Lượt 1} & \textbf{Lượt 2} & \textbf{Lượt 3 (mới nhất)} \\
\midrule
Phần cứng          & CPU     & CPU     & CPU + GPU \\
Số agent lúc huấn luyện & 12--16  & 48      & \textbf{100 (đúng mức tối đa bài báo)} \\
Độ dài một lượt mô phỏng & 300 giây & 420 giây & \textbf{600 giây (đúng bài báo)} \\
Số episode huấn luyện   & 800--1.800 & 2.000--2.800 & 3.000--4.000 \\
\bottomrule
\end{tabular}
\end{table}

\subsection{Các kiểm chứng bổ sung (sau lượt 3)}

\begin{itemize}[leftmargin=1.2em]
  \item \textbf{Huấn luyện lặp lại với 3 hạt giống ngẫu nhiên (seed) độc lập} cho cả 3 chính
        sách, để kiểm tra kết quả có ổn định hay chỉ là một lần chạy may mắn. Kết quả: kết luận
        chính vẫn giữ vững qua cả 3 lần chạy độc lập.
  \item \textbf{Quét toàn bộ hệ số pha (phasing factor) của chòm Walker-Delta} (48 giá trị khả
        dĩ) để tìm xem có giá trị nào cho đúng con số bài báo công bố hay không. Kết quả: không
        có giá trị nào khớp — đây là kết quả âm tính đáng tin vì đã thử hết toàn bộ khả năng.
  \item \textbf{Bổ sung so sánh với chòm vệ tinh OneWeb} (có trong bảng tham số của bài báo
        nhưng chưa từng được đưa vào so sánh ở 2 lượt đầu).
\end{itemize}

\subsection{Hai thử nghiệm bổ sung tiếp theo}

\begin{itemize}[leftmargin=1.2em]
  \item \textbf{Huấn luyện riêng (native) cho từng chòm vệ tinh}, thay vì chỉ chuyển giao
        chính sách đã huấn luyện trên chòm chính sang chòm khác (zero-shot) như trước. Kết
        quả (Bảng~\ref{tab:native}): huấn luyện riêng cho số chuyển giao thất bại thấp hơn
        hẳn ở mức tải trung bình--cao, đánh đổi bằng thông lượng thấp hơn một chút.
  \item \textbf{Kiểm tra độ nhạy với kích thước lô dữ liệu huấn luyện (batch size)} — thử
        `batch=4' thay cho `batch=8' của lượt 3 (giá trị `batch=8' vốn đã thấp hơn 32 mà bài
        báo công bố, do giới hạn bộ nhớ GPU). Kết quả: \textbf{không có khác biệt đáng kể} —
        mọi số liệu của `batch=4' nằm gọn trong khoảng dao động tự nhiên đã đo được giữa
        3 lần huấn luyện độc lập với `batch=8'. Kích thước lô, ít nhất trong khoảng đã thử,
        không phải yếu tố ảnh hưởng tới kết quả.
\end{itemize}

\begin{table}[h]
\centering\small
\caption{Huấn luyện riêng (native) so với chuyển giao chính sách (zero-shot)}
\label{tab:native}
\begin{tabular}{lcc}
\toprule
 & \textbf{Phase 1-a} & \textbf{Hybrid} \\
\midrule
Số chuyển giao thất bại thấp hơn & 2,6--2,7 lần (tải 20--40 UE) & 2,4--5,1 lần (tải 20--70 UE) \\
Đánh đổi thông lượng & thấp hơn 8--13\% & thấp hơn 3--12\% \\
\bottomrule
\end{tabular}
\end{table}

\section{Kết quả chính}

Kết quả nổi bật nhất: khi tăng dần quy mô huấn luyện qua 3 lượt (đúng số agent, đúng độ dài
mô phỏng bài báo công bố), khoảng cách hiệu năng giữa ILCHO và các phương pháp so sánh tốt
nhất thu hẹp rất nhanh (Bảng~\ref{tab:hof}), cho thấy \textbf{ngân sách huấn luyện chính là yếu
tố chính} khiến lượt tái hiện đầu chưa khớp bài báo, không phải do sai phương pháp.

\begin{table}[h]
\centering
\caption{Số lần chuyển giao thất bại (HOF) trung bình mỗi thiết bị, khi hệ thống có 100 người
dùng cùng lúc — càng thấp càng tốt}
\label{tab:hof}
\begin{tabular}{lccc}
\toprule
\textbf{Phương pháp} & \textbf{Lượt 1} & \textbf{Lượt 2} & \textbf{Lượt 3} \\
\midrule
ILCHO (đề xuất)        & $\approx$47,4 & 21,1 & \textbf{4,95} (4,52 khi lấy trung bình 3 seed) \\
So với phương pháp tốt nhất khác & --   & 6,3 lần & \textbf{1,34 lần} \\
\bottomrule
\end{tabular}
\end{table}

Các kết luận khác đã tái hiện được, vững chắc qua khảo sát độ nhạy tham số:

\begin{itemize}[leftmargin=1.2em]
  \item ILCHO giữ số lần chuyển giao thất bại thấp hơn hẳn 2 phương pháp truyền thống
        (khoảng cách gần nhất, thời gian nhìn thấy lâu nhất) ở \textbf{mọi} mức số người dùng
        đã thử (5 đến 100).
  \item Độ công bằng giữa các người dùng (Jain fairness index) đạt 0,995--0,998, khớp/tốt hơn
        con số bài báo công bố (0,9798).
  \item Hình học quỹ đạo và mô hình liên kết vô tuyến khớp sát số liệu bài báo (sai lệch vài
        phần trăm).
\end{itemize}

Hai tuyên bố của bài báo \textbf{không tái hiện được}, đã kiểm chứng nhất quán qua cả 3 lượt
nên đây là phát hiện thật, không phải do thiếu sót:

\begin{itemize}[leftmargin=1.2em]
  \item Bài báo cho rằng ILCHO có thông lượng cao nhất trong mọi phương pháp — trong bản tái
        hiện này, phương pháp HSNF (tham lam tối đa hoá thông lượng tức thời) luôn cao hơn.
  \item Bài báo cho rằng hàm thưởng tuyến tính (một biến thể được thử nghiệm) huấn luyện bất
        ổn — trong bản tái hiện này, nó hội tụ ổn định, gần tương đương hàm thưởng chính.
\end{itemize}

\section{Hạn chế và những gì chưa làm được}

\begin{itemize}[leftmargin=1.2em]
  \item \textbf{Một số tham số bài báo không công bố} (tần số sóng mang chính xác, ngưỡng
        thực thi chuyển giao, các hằng số hàm thưởng, hệ số pha chòm vệ tinh, siêu tham số
        gốc của 2 phương pháp so sánh). Đã tự chọn giá trị hợp lý và kiểm chứng độ nhạy với
        2 trong 5 tham số này (kết luận không đổi dù giá trị là bao nhiêu) — 3 tham số còn lại
        chưa kiểm chứng được đầy đủ. Đây là giới hạn \textbf{không thể tự đóng được} nếu không
        liên hệ trực tiếp nhóm tác giả bài báo.
  \item Hai phương pháp so sánh (HSNF, LBSH) đang dùng bản triển khai rút gọn của mình, chưa
        cài đúng theo 2 bài báo tham chiếu gốc — đây là việc còn có thể làm thêm nếu có thời
        gian.
  \item Do giới hạn bộ nhớ GPU hiện có, kích thước lô dữ liệu huấn luyện (`batch=8', thử thêm
        `batch=4') vẫn thấp hơn `batch=32' mà bài báo công bố — dù đã kiểm chứng không có
        khác biệt đáng kể giữa 4 và 8, chưa thể loại trừ hoàn toàn khả năng 32 khác biệt.
  \item Mới kiểm chứng ổn định qua 3 lần huấn luyện độc lập (khuyến nghị thống kê tối thiểu);
        có thể tăng lên 5 lần nếu muốn kết quả chặt hơn.
\end{itemize}

\section{Kế hoạch tiếp theo}

\begin{enumerate}[leftmargin=1.2em]
  \item Đọc và cài đặt lại đúng 2 phương pháp so sánh (HSNF, LBSH) theo đúng 2 bài báo tham
        chiếu gốc — đây là việc có khả năng làm rõ nhất liệu 2 kết luận chưa tái hiện được ở
        Mục 4 là do bản triển khai rút gọn hay là khác biệt thật.
  \item Cân nhắc liên hệ nhóm tác giả bài báo để xin các tham số chưa công bố.
  \item Nếu còn thời gian: tăng số lần huấn luyện độc lập từ 3 lên 5 để khoảng tin cậy chặt
        hơn.
\end{enumerate}

\section{Kết luận}

Sau 3 lượt tái hiện, tuyên bố cốt lõi của bài báo — một cơ chế chuyển giao điều kiện điều
khiển bằng học tăng cường đa tác tử giữ số chuyển giao thất bại thấp và ổn định trong chòm vệ
tinh mật độ cao, nơi các phương pháp truyền thống thất bại khi tải tăng — \textbf{đã tái hiện
được rõ ràng và vững chắc}. Khoảng cách với bài báo ở vùng tải cao nhất đã thu hẹp mạnh nhờ
huấn luyện đúng quy mô bài báo công bố. Một số chi tiết (thông lượng cao nhất, độ ổn định của
một biến thể hàm thưởng) chưa tái hiện được và đã được báo cáo trung thực như các phát hiện
riêng, không che giấu.

\vspace{1em}
\noindent\emph{Toàn bộ mã nguồn, số liệu chi tiết và các báo cáo kỹ thuật đầy đủ (tiếng Việt)
được lưu tại kho mã nguồn của đề tài.}

\end{document}
